\section{The scope boundary (constructive)}

Two systems are constructed with identical local reactant wells (same \(M\), \(C\), \(K\)) but different global barriers. ChemSA returns the same class and identical poles for both. Because the two systems share the well, friction, and barrier frequency, the Kramers--Grote--Hynes prefactor cancels between them exactly, leaving \(k_A/k_B = e^{(\Delta G_B - \Delta G_A)/k_BT} = e^{5} \approx 148\): no absolute Kramers rate is evaluated, only this prefactor-free ratio, since the friction used is deeply underdamped and outside the moderate-to-high-friction regime the prefactor itself is derived for. Identical local stability architecture, rates differing by 148-fold: local classification cannot determine rate. ChemSA therefore classifies mode stability and does not infer reaction rate from local architecture alone. A transition state (a stiffness matrix with a negative eigenvalue) is always classified unstable, never as a stable-well mode, for all frictions in the scalar inverted-harmonic model tested.

\subsection*{Where the coordinate does not exist}

The qualifying rule of Section~1.3 restricts \(\chi\) to a stable second-order response. There is a sharper statement available on the other side of the surface, and it is algebraic rather than evidential. For a stable mode the characteristic discriminant is \(\gamma^2/4 - \omega_0^2\), which vanishes at \(\gamma = 2\omega_0\); that vanishing is the boundary \(\chi = 1\). For an inverted mode, a transition state with negative curvature along the reaction coordinate,
\[
\ddot q + \gamma \dot q - \omega_b^2 q = 0, \qquad \text{discriminant} = \frac{\gamma^2}{4} + \omega_b^2 ,
\]
and this cannot vanish for real \(\gamma\) and nonzero \(\omega_b\). The roots remain real and distinct at every friction, one of them always positive. A dimensionless friction remains perfectly well defined at a barrier. Writing it \(\alpha_b = \gamma/2\omega_b\), the normalized barrier friction, it parameterizes \(\lambda_+\) and \(\kappa_b\) in the scalar passive model and governs how much of an incoming flux is transmitted. It is deliberately not called a mechanical damping ratio, because that name carries a critical value and this quantity has none. What is absent is the boundary: the roots never coalesce, so nothing changes character at \(\alpha_b = 1\) and the coordinate carries no landmark on the side of the surface where reactions occur.

This is the reason a damping ratio cannot carry a rate, stated without a counterexample. The quantity that does govern transmission through a barrier is the positive reactive pole \(\lambda_+ = -\gamma/2 + \sqrt{\gamma^2/4 + \omega_b^2}\), or equivalently the transmission factor \(\kappa_b = \lambda_+/\omega_b\) of Kramers and of Grote and Hynes \citep{hanggi1990}, which is monotonically decreasing and so has no maximum either. Locating exceptional points in parametric quadratic eigenproblems is itself established: continuation methods and the open-source \texttt{EasterEig} implementation compute defective eigenvalues and their Puiseux expansions directly \citep{eastereig}. The contribution here is not the location of such points but the conservative workflow around them, namely classification of a supplied representation, refusal of neighbourhoods that finite precision cannot resolve, enforcement of representation provenance, and channel-declared response geometry. The engine already enforces this: a stiffness matrix with a negative eigenvalue is classified as a saddle and \(\chi\) is withheld for the reactive coordinate at every friction tested, with the reactive pole reported in its place. Stable transverse modes at the same saddle retain their own damping ratios, which is why the statement is about the reactive direction rather than about the whole generator. Even that pole is only a prefactor, since the activation exponent is not a property of the local generator, which is the same boundary the constructive counterexample of Section~3 establishes by a different route.
